\section{Methods}
\label{sec:methods}
We treated inference of the TSC as a supervised learning problem, with radio and X-ray observations as the input data and the pseudo-TSC label defined in section \ref{sec:sims}.
Note that the pseudo-TSC is not a directly-measured time since pericentric passage; it instead measures the elapsed time since the merger score crosses a particular threshold.
The generation procedure for the radio and X-ray mocks from simulated data is described in sections \ref{sec:sims} and \ref{sec:fwd}.

We primarily trained a convolutional neural network (CNN), which has shown strong performance on similar astrophysical and cosmological inference problems.
The CNN independently processed each of the three spatial projections of an input cluster using a shared encoder comprising four convolutional blocks, with 32, 64, 128, and 256 output channels.
Each block of the CNN has two convolutional layers, with batch normalization and ReLU (Rectified Linear Unit) nonlinear activations, followed by spatial downsampling.
We then used a global average pooling layer to calculate a 256-dimensional representation of each cluster projection.
The three projection representations are averaged, which makes the prediction independent of projection ordering.
This vector is then passed to a regression head, constructed from a 64-hidden-unit layer with dropout.
In the joint radio and X-ray model, each modality has a separate encoder, with the encodings concatenated before the regression step.

While training, we augmented our set of radio and X-ray projections with random reflections, and $90^\circ$, $180^\circ$, and $270^\circ$ rotations, applied consistently to both modalities.
\todo{clean map experiments}
Observationally forward-modelled data was preprocessed using the procedure described in section \ref{sec:fwd}.
We used AdamW to optimize the the CNN for 60 epochs, with an initial learning rate of $3\times 10^{-4}$, weight decay of $10^{-3}$, a cosine learning-rate schedule, and a batch size of 32.
We used a Huber loss \todo{Huber loss formula} with transition scale $0.5$ Gyr to reduce the influence of large residuals.

To evaluate performance, we used five-fold cross validation, with data from each cluster (projections, noise realizations, and available snapshots) always being in the same fold to prevent data leakage.
Within each outer training partition, 15\% of clusters were reserved for training checkpoint selection; we retained the checkpoint with the the lowest inner-validation squared error, and excluded the outer test partition from checkpoint selection.
Repeating this procedure for each of the five outer folds produced an out-of-fold (OOF) prediction for each cluster.
Overall, 68\% of clusters were used for training, 12\% were used for inner validation, and 20\% were used for OOF predictive testing.

We measured performance using the coefficient of determination, $r^2$, with overall $r^2$ calculated from the combined OOF predictions, rather than by averaging the $r^2$ values of individual folds.
Our primary experiments were repeated with five random seeds, with comparisons between conditions using shared seeds to reduce noise in comparisons.
We measured scatter across seeds to estimate variation between runs, but this cannot be interpreted as a calibrated uncertainty interval for the TSC.

To evaluate the increased predictive power from access to raw image data, we compared our CNN model with several scalar prediction baselines.
A five-feature ridge regression model was fit from log(halo mass), log(halo radius), log(integrated emission), mean intensity in the arcsinh-stretched image, and the fraction of pixels with non-zero intensity.
We standardized features over the training data, with ridge regularization strength selected within the training sample.
These comparisons tested whether the raw image data provided extra predictive power over summary statistics of a cluster's scale, emission amplitude, and simple properties of the intensity distribution.

We also compared the CNN performance to explicit image morphology descriptions.
For radio maps, we constructed thresholded excursion sets at 12 intensity quantiles, forming a 48-dimensional feature vector containing their perimeter, Euler characteristic, and numbers of connected components and holes.
This was supplied to an XGBoost regressor.
All projections of a cluster remained grouped during cross-validation, with their predictions averaged for cross-cluster comparison.
For X-ray data, we fit a ridge-regression model on aperture-based concentration, centroid variation, ellipticity, and integrated emission.
These were defined relative to the image extent, also with features averaged across projections.
\todo{should probably define some of these more clearly}

Several complementary analyses considered the role of mass and emission amplitude.
To test the contribution of explicit amplitude information, one analysis ran shape-normalization \todo{more detail}, with each projection then rescaled to a common integrated emission before applying the intensity arcsinh-stretch.
Note that this normalization did not remove correlations between mass, morphology, and luminosity.
We also considered predictive performance within mass bins, and calculated partial-rank correlations between predictions and labels after accounting for halo mass.

For analyses in which we conditioned on cluster mass, we concatenated $\log_{10} (M_\text{500c})$ to the image representation before being fed into the regression head.
To approximate uncertainty associated with weak-lensing masses, we perturbed simulated log-masses with Gaussian noise with widths sampled from the empirical distribution of weak-lensing fractional mass errors.
We used the available fractional uncertainties on $M_\text{200c}$ estimates to approximate uncertainties in the converted $M_\text{500c}$ values.

The experiments that have access to all three cluster projections provide a benchmark in which several views of the same cluster are available, which can be compared with the model performance on the realistic case individual observed images from a single viewing direction.
The latter can be interpreted alongside the corresponding observationally forward-modelled tests.
Because the observed clusters lack independent pseudo-TSC labels, agreement between radio and X-ray predictions provides a consistency check rather than a direct accuracy measurement.